\documentclass[11pt,a4paper]{article}

\usepackage[utf8]{inputenc}
\usepackage[T1]{fontenc}
\usepackage[margin=2.5cm]{geometry}
\usepackage{amsmath,amssymb,bm}
\usepackage{booktabs}
\usepackage{array}
\usepackage{enumerate}
\usepackage{xcolor}
\usepackage{hyperref}
\usepackage{tikz}

\hypersetup{colorlinks=true, linkcolor=blue, urlcolor=blue}

\newcommand{\points}[1]{\textbf{(#1\,\%)}}
\newcommand{\subpoints}[1]{\textit{(#1\,\%)}}
\newcommand{\sol}{\noindent\textbf{Solution}}
% \boxed is already provided by amsmath

\title{TMA4268 V2026 Mock Exam 2\\[0.3em]
\large Solution Proposal}
\author{Compiled for Anders Bekkevard\\[0.2em]
\small Companion to \texttt{mock-exam-2.tex} (same directory).}
\date{Mock for: May 18, 2026}

\begin{document}
\maketitle

\noindent\textit{This document is a worked-solution proposal in the style of the official Stefanie/Sara solutions
to the 2024 and 2025 finals. Point values are quoted in the heading of each solved sub-part. Partial-credit
hints are inline. Refer back to \texttt{mock-exam-2.tex} for the problem statements.}

\vspace{0.6em}
\hrule
\vspace{1em}

%=========================================================
\section*{Problem 1 \points{10} --- Fill-in-the-blank concepts}
%=========================================================

\sol{} \emph{(1\,P per blank.)}

\begin{enumerate}[(1)]
  \item \textbf{validation}
  \item \textbf{cross-validation}
  \item \textbf{bootstrap}
  \item \textbf{forward stepwise}
  \item \textbf{lasso}
  \item \textbf{sparse}
  \item \textbf{principal component analysis}
  \item \textbf{local}
  \item \textbf{dendrogram}
  \item \textbf{boosting}
\end{enumerate}

\noindent\emph{Grading: 1\,P per correct blank. ``Holdout'' is acceptable for (1) only if accompanied by
``validation''; the passage contrasts validation against the test set. ``Bagging'' for (3) is wrong:
bagging is an \emph{ensemble} technique built on top of the bootstrap, but the resampling tool itself
is the bootstrap.}

\vspace{0.6em}
\hrule
\vspace{1em}

%=========================================================
\section*{Problem 2 \points{28} --- Multiple choice, T/F, short numeric}
%=========================================================

\subsection*{a) The bootstrap \subpoints{3}}
\sol{} \emph{(0.75\,P per statement.)}
\begin{enumerate}[(i)]
  \item \textbf{True}. The defining feature of the bootstrap is sampling \emph{with} replacement, same size $n$.
  \item \textbf{False}. A single resample gives one realisation of $\hat\theta^*$; the SE estimate is the
        standard deviation of $\hat\theta^{*1},\dots,\hat\theta^{*B}$, which requires $B$ resamples (typically
        $B\geq 200$ for SEs, $\geq 1000$ for percentile CIs).
  \item \textbf{True}. $P(\text{obs } i \notin \text{sample}) = (1-1/n)^n \to 1/e \approx 0.368$.
  \item \textbf{True}. The \emph{percentile interval} is one of the two standard bootstrap-CI constructions;
        the other is the normal-approximation interval $\hat\theta \pm z_{\alpha/2}\cdot\widehat{\mathrm{SE}}_{\text{boot}}$.
        The percentile method makes no Gaussianity assumption and is directly read off the bootstrap
        distribution's quantiles, hence its appeal for statistics with skewed or otherwise non-normal
        sampling distributions.
\end{enumerate}

\subsection*{b) Cross-validation and its pitfalls \subpoints{3}}
\sol{} \emph{(0.75\,P per statement.)}
\begin{enumerate}[(i)]
  \item \textbf{True}. LOOCV averages $n$ highly-correlated training fits (they share $n-1$ observations), so
        its variance as a test-error estimator is high. $k$-fold with $k=5$ or $10$ has more independent folds
        and lower variance, paid for by a small upward bias (each fit uses only $(k-1)n/k$ observations).
  \item \textbf{False}. Pre-selecting predictors on the \emph{full} training set lets information from each
        fold's held-out observations leak into the predictor screen. The CV error is then \emph{optimistically
        biased}. The fix is to put the predictor screening \emph{inside} each CV fold.
  \item \textbf{True}. The validity of $k$-fold CV rests on the held-out fold being exchangeable with the
        training fold. Time-series data violate this (autocorrelation, drift), and naive $k$-fold mixes future
        and past observations. Time-series CV (rolling/expanding window) is the remedy.
  \item \textbf{False}. The roles are swapped: the \emph{inner} folds are used for hyperparameter tuning, and
        the \emph{outer} folds are used for unbiased performance assessment of the whole tuning procedure.
\end{enumerate}

\subsection*{c) Subset selection \subpoints{3}}

\sol{} \subpoints{1} \textbf{(i)} Best-subset on $p=8$ predictors fits every subset, including the null:
$$2^p = 2^8 = \boxed{256} \text{ candidate models.}$$

\medskip
\sol{} \subpoints{1} \textbf{(ii)} Forward stepwise fits the null model plus, at each step
$k=0,1,\dots,p-1$, the $p-k$ candidates that add one of the remaining predictors. Total:
$$1 + \sum_{k=0}^{p-1}(p-k) = 1 + p + (p-1) + \cdots + 1 = 1 + \frac{p(p+1)}{2}
   = 1 + \frac{8\cdot 9}{2} = \boxed{37} \text{ candidate models.}$$

\medskip
\sol{} \subpoints{1} \textbf{(iii)} \textbf{False.} Backward stepwise starts from the \emph{full} model with all
$p$ predictors and removes one at a time. When $p>n$ the full OLS fit is not even defined (the design matrix
has more columns than rows), so backward stepwise cannot be applied in the high-dimensional regime. Forward
stepwise, by contrast, can run until the model size hits $\min(n-1,p)$.

\subsection*{d) Cubic and natural splines: counting degrees of freedom \subpoints{3}}

\sol{} \subpoints{1} \textbf{(i)} A cubic regression spline with $K$ interior knots, fit on top of a global
intercept, uses $K+3$ basis functions for the spline term (the truncated-power basis: $x,x^2,x^3$ plus one
truncated cube per knot). With $K=5$:
$$K + 3 = 5 + 3 = \boxed{8} \text{ coefficients (excluding the intercept).}$$

\medskip
\sol{} \subpoints{1} \textbf{(ii)} A \emph{natural} cubic spline imposes linearity beyond the boundary knots,
which removes two effective parameters from each end (four constraints, of which two are absorbed by the
already-required smoothness). The standard count is $K+1$ basis functions (excluding the intercept):
$$K + 1 = 5 + 1 = \boxed{6} \text{ coefficients (excluding the intercept).}$$

\medskip
\sol{} \subpoints{1} \textbf{(iii)} \textbf{True.} A cubic spline is, by construction, $C^2$ at every knot:
continuous in value, first derivative, and second derivative. Only the third derivative is permitted to jump,
which is what allows the piecewise cubic to bend differently in adjacent intervals.

\subsection*{e) Logistic regression: odds, log-odds, probability \subpoints{3}}

\sol{} \subpoints{1} \textbf{(i)} A change of $\Delta x = 10$ years in $\texttt{age}$ multiplies the odds by
$$e^{\hat\beta_{\text{age}}\cdot 10} = e^{0.04\cdot 10} = e^{0.4} \approx \boxed{1.49}.$$
So the older individual's odds are about $1.49\times$ those of the younger.

\medskip
\sol{} \subpoints{1} \textbf{(ii)} Using the logistic transform,
$$\hat p = \frac{e^{\hat\eta}}{1 + e^{\hat\eta}} = \frac{e^{-0.5}}{1 + e^{-0.5}}
       = \frac{0.6065}{1.6065} \approx \boxed{0.38}.$$

\medskip
\sol{} \subpoints{1} \textbf{(iii)} \textbf{False.} $\beta_j$ is a change in \emph{log-odds}, not probability.
The change in probability for a unit increase in $x_j$ depends nonlinearly on the current value of $\hat p$
(largest near $\hat p=0.5$, smallest near $0$ or $1$). The clean linear interpretation is on the log-odds /
odds-multiplication scale, not on the probability scale.

\subsection*{f) Confusion matrix interpretation \subpoints{4}}

We have $n=1500$ patients, with $300$ true positives in total (TP $=240$, FN $=60$) and $1200$ true negatives
in total (TN $=1080$, FP $=120$).

\medskip
\sol{} \subpoints{1} \textbf{(i)} Accuracy $= (\text{TP}+\text{TN})/n
   = (240 + 1080)/1500 = 1320/1500 = \boxed{0.88}.$

\medskip
\sol{} \subpoints{1} \textbf{(ii)} Sensitivity $= \text{TP}/(\text{TP}+\text{FN})
   = 240/300 = \boxed{0.80}.$

\medskip
\sol{} \subpoints{1} \textbf{(iii)} Specificity $= \text{TN}/(\text{TN}+\text{FP})
   = 1080/1200 = \boxed{0.90}.$

\medskip
\sol{} \subpoints{1} \textbf{(iv)} Precision $= \text{TP}/(\text{TP}+\text{FP})
   = 240/(240+120) = 240/360 = \boxed{0.67}.$

\subsection*{g) $K$-means clustering \subpoints{3}}
\sol{} \emph{(0.75\,P per statement.)}
\begin{enumerate}[(i)]
  \item \textbf{False}. $K$-means descends a non-convex objective; for a given initialization it converges
        to a \emph{local} optimum. Different random starts can land in different basins; the standard
        remedy is to restart $K$-means many times and keep the best.
  \item \textbf{True}. Each step (reassign-to-nearest-centroid, recompute-centroid-as-mean) is guaranteed
        to weakly decrease the within-cluster sum of squares, which is what gives the algorithm its
        convergence guarantee (to a local minimum).
  \item \textbf{True}. Without standardization, large-scale variables dominate the squared Euclidean
        distance and the clusters are essentially partitioned along those directions. Standardizing
        ($z$-scores) is the usual fix.
  \item \textbf{False}. $K$-means is not nested across $K$. Increasing $K$ from $4$ to $5$ can reorganise
        the partition globally; the only structure that is monotone in $K$ is hierarchical clustering
        (with single/complete linkage).
\end{enumerate}

\subsection*{h) Tree ensembles: bagging, random forests, boosting \subpoints{4}}
\sol{} \emph{(1\,P per statement.)}
\begin{enumerate}[(i)]
  \item \textbf{True}. Smaller $m$ forces each split to consider a more restricted predictor set, decorrelating
        the trees. Per the bagging variance formula
        $\mathrm{Var}(\bar f)=\rho\sigma^2+(1-\rho)\sigma^2/B$, lowering $\rho$ lowers the floor. (See P3a.)
  \item \textbf{False}. Bagging and random forests do \emph{not} overfit in $B$: as $B\to\infty$ the variance
        of the averaged predictor monotonically decreases (toward $\rho\sigma^2$). $B$ is chosen large enough
        that the OOB or test error stabilizes, not by CV. (Contrast with boosting, where $B$ \emph{is} a real
        tuning parameter.)
  \item \textbf{True}. For squared loss $L(y,F)=\tfrac12(y-F)^2$, the negative functional gradient with
        respect to $F$ is $-\partial L/\partial F = y - F$, i.e.\ the current residual. So gradient boosting
        with squared loss is residual-fitting.
  \item \textbf{True}. The textbook defaults are $m=\lfloor\sqrt p\rfloor$ for classification and
        $m=\lfloor p/3\rfloor$ for regression.
\end{enumerate}

\subsection*{i) Direction of effect \subpoints{2}}
\sol{} \emph{(1\,P per statement.)}
\begin{enumerate}[(i)]
  \item \textbf{True}. The smoothing-spline effective degrees of freedom $\mathrm{tr}(S_\lambda)$ shrink from
        $n$ (interpolation, $\lambda=0$) toward $2$ (OLS line, $\lambda\to\infty$). Increasing $\lambda$ gives
        a smoother, less-wiggly fit.
  \item \textbf{False}. The direction is the opposite. \emph{Smaller} $\nu$ means each tree contributes \emph{less}
        to the running ensemble, so \emph{more} trees $B$ are needed to fit the data well. The textbook rule of
        thumb is that halving $\nu$ roughly doubles the required $B$.
\end{enumerate}

\vspace{0.6em}
\hrule
\vspace{1em}

%=========================================================
\section*{Problem 3 \points{16} --- Theory and hand calculations}
%=========================================================

\subsection*{a) The mathy one: variance of an average of correlated predictors \subpoints{8}}

\sol{} \subpoints{4} \textbf{(i)} Let $X_b = \hat f^{*b}(x_0)$ for $b=1,\dots,B$ to shorten notation. By
hypothesis, $\mathrm{Var}(X_b)=\sigma^2$ and $\mathrm{Cov}(X_b,X_{b'})=\rho\sigma^2$ for $b\neq b'$.
Using linearity of variance and the bilinear expansion of $\mathrm{Var}(\sum_b X_b)$:
\begin{align*}
\mathrm{Var}\!\left(\bar f_{\text{bag}}(x_0)\right)
   &= \mathrm{Var}\!\left(\frac{1}{B}\sum_{b=1}^B X_b\right)
   = \frac{1}{B^2}\,\mathrm{Var}\!\left(\sum_{b=1}^B X_b\right) \\[2pt]
   &= \frac{1}{B^2}\!\left[\sum_{b=1}^B\mathrm{Var}(X_b) + \sum_{b\neq b'}\mathrm{Cov}(X_b,X_{b'})\right] \\[2pt]
   &= \frac{1}{B^2}\!\left[B\,\sigma^2 + B(B-1)\,\rho\,\sigma^2\right]
      \qquad\text{(there are $B(B-1)$ ordered off-diagonal pairs)} \\[2pt]
   &= \frac{\sigma^2}{B} + \frac{B-1}{B}\,\rho\,\sigma^2 \\[2pt]
   &= \frac{1}{B}\sigma^2 + \left(1 - \frac{1}{B}\right)\rho\sigma^2 \\[2pt]
   &= \rho\sigma^2 + \frac{1}{B}\sigma^2 - \frac{1}{B}\rho\sigma^2 \\[2pt]
   &= \boxed{\;\rho\,\sigma^2 \;+\; \frac{1-\rho}{B}\,\sigma^2.\;}
\end{align*}

\noindent\emph{Grading: 1\,P for splitting $\mathrm{Var}(\sum X_b)$ into diagonal + off-diagonal terms;
1\,P for the correct count $B$ diagonal and $B(B-1)$ off-diagonal; 1\,P for simplifying to
$\sigma^2/B + (B-1)\rho\sigma^2/B$; 1\,P for the final form $\rho\sigma^2+(1-\rho)\sigma^2/B$. Deduct
1 per algebra slip.}

\medskip
\sol{} \subpoints{1} \textbf{(ii)} Letting $B\to\infty$ with $\rho,\sigma^2$ fixed:
$$\lim_{B\to\infty}\mathrm{Var}\!\left(\bar f_{\text{bag}}(x_0)\right) = \boxed{\rho\,\sigma^2}.$$
The $(1-\rho)\sigma^2/B$ term vanishes, leaving the floor $\rho\sigma^2$ that cannot be removed by adding
more bootstrap-trained predictors.

\medskip
\sol{} \subpoints{2} \textbf{(iii)} For bagged trees built on $B$ bootstrap samples of the \emph{same} training
data, the trees tend to be highly correlated: a strong predictor dominates the top splits in essentially
every tree, so the $B$ trees look alike. Then $\rho$ is close to $1$ and the limit $\rho\sigma^2$ in (ii)
is close to $\sigma^2$ --- averaging more trees barely helps.

The \emph{random forest} algorithm attacks this floor by, at each split, restricting the candidate predictors
to a random subset of size $m<p$. This deliberately stops the dominant predictor from being used at every
top split and \emph{decorrelates} the trees: $\rho$ falls, and so does the asymptotic floor $\rho\sigma^2$.
The price is a small per-tree bias increase, paid for by the variance reduction.

\medskip
\sol{} \subpoints{1} \textbf{(iv)} The extreme case is $\rho=0$ (i.i.d.\ predictions). The formula then
collapses to
$$\mathrm{Var}\!\left(\bar f_{\text{bag}}\right) = 0\cdot\sigma^2 + \frac{1-0}{B}\sigma^2 = \boxed{\sigma^2/B},$$
which is the familiar variance of the mean of $B$ i.i.d.\ random variables.

\subsection*{b) Hierarchical clustering by hand: single linkage \subpoints{5}}

\sol{} \subpoints{3} \textbf{(i)} Reading off the upper triangle of $D$:
$$d_{12}=2,\;d_{13}=6,\;d_{14}=10,\;d_{15}=9,\;d_{23}=5,\;d_{24}=9,\;d_{25}=8,\;d_{34}=4,\;d_{35}=5,\;d_{45}=3.$$

\paragraph{Merge 1 (height 2).} Smallest dissimilarity is $d_{12}=2$. Merge $\{1\}\cup\{2\}\to\{1,2\}$
at height $\mathbf{h_1=2}$. Under \emph{single} linkage the new row/column is the \emph{minimum} of the merged rows:
\begin{align*}
d(\{1,2\},3) &= \min(d_{13},d_{23}) = \min(6,5) = 5, \\
d(\{1,2\},4) &= \min(d_{14},d_{24}) = \min(10,9) = 9, \\
d(\{1,2\},5) &= \min(d_{15},d_{25}) = \min(9,8) = 8.
\end{align*}
The recomputed $4\times 4$ dissimilarity matrix is
$$D^{(1)} = \bordermatrix{
       & \{1,2\} & 3 & 4 & 5 \cr
\{1,2\}& 0       & 5 & 9 & 8 \cr
3      & 5       & 0 & 4 & 5 \cr
4      & 9       & 4 & 0 & 3 \cr
5      & 8       & 5 & 3 & 0 }.$$

\paragraph{Merge 2 (height 3).} Smallest off-diagonal of $D^{(1)}$ is $d_{45}=3$. Merge $\{4\}\cup\{5\}\to\{4,5\}$
at height $\mathbf{h_2=3}$. With single linkage:
\begin{align*}
d(\{4,5\},\{1,2\}) &= \min(d_{14},d_{24},d_{15},d_{25}) = \min(10,9,9,8) = 8, \\
d(\{4,5\},3)       &= \min(d_{34},d_{35}) = \min(4,5) = 4.
\end{align*}
The recomputed $3\times 3$ dissimilarity matrix is
$$D^{(2)} = \bordermatrix{
       & \{1,2\} & 3 & \{4,5\} \cr
\{1,2\}& 0       & 5 & 8       \cr
3      & 5       & 0 & 4       \cr
\{4,5\}& 8       & 4 & 0       }.$$

\paragraph{Merge 3 (height 4).} Smallest off-diagonal of $D^{(2)}$ is $d(\{4,5\},3)=4$. Merge
$\{3\}\cup\{4,5\}\to\{3,4,5\}$ at height $\mathbf{h_3=4}$. Single-linkage update for the only remaining pair:
$$d(\{1,2\},\{3,4,5\}) = \min(d_{13},d_{14},d_{15},d_{23},d_{24},d_{25}) = \min(6,10,9,5,9,8) = 5.$$

\paragraph{Merge 4 (height 5).} The final merge of $\{1,2\}$ with $\{3,4,5\}$ occurs at height
$\mathbf{h_4=5}$.

\noindent\emph{Grading: 1\,P for merge 1 with the correctly updated $D^{(1)}$; 1\,P for merge 2 with
correctly updated $D^{(2)}$; 1\,P for both remaining heights $h_3=4$ and $h_4=5$. Deduct 0.5 per
arithmetic slip on the linkage minimum.}

\medskip
\sol{} \subpoints{1} \textbf{(ii)} Dendrogram with the four fusion heights $2,3,4,5$ on the $y$-axis:
\begin{center}
\begin{tikzpicture}[scale=0.85]
  % axis
  \draw[->] (-0.5,0) -- (-0.5,6) node[above]{$h$};
  \foreach \y in {0,1,2,3,4,5} \draw (-0.55,\y) -- (-0.45,\y) node[left,xshift=-3pt]{\y};
  % leaves at y=0, ordering 1,2,3,4,5
  \foreach \x/\lab in {0/1, 1/2, 2.5/3, 4/4, 5/5} \node at (\x,-0.3) {$\lab$};
  % merge 1 & 2 at h=2
  \draw (0,0) -- (0,2) -- (1,2) -- (1,0);
  % merge 4 & 5 at h=3
  \draw (4,0) -- (4,3) -- (5,3) -- (5,0);
  % merge 3 with {4,5} at h=4
  \draw (2.5,0) -- (2.5,4) -- (4.5,4) -- (4.5,3);
  % top merge at h=5
  \draw (0.5,2) -- (0.5,5) -- (3.5,5) -- (3.5,4);
\end{tikzpicture}
\end{center}

\medskip
\sol{} \subpoints{1} \textbf{(iii)} \textbf{No, the first merge would not change.} The first merge only
compares pairs of singletons, and the smallest entry of $D$ ($d_{12}=2$) is the smallest single-pair
dissimilarity regardless of linkage. Linkage rules only start to matter from the second merge onwards,
when at least one cluster has size $\geq 2$ and the rule decides how to summarise it.

\subsection*{c) The bootstrap, by hand \subpoints{3}}

\sol{} \subpoints{1} \textbf{(i)} Each draw in the bootstrap sample is uniform on $\{1,\dots,n\}$ with
replacement, so the probability that a single draw is \emph{not} observation $i$ is $1-1/n$. The $n$ draws
are independent, so the probability that observation $i$ is missed by \emph{all} $n$ draws is
$$P(i\notin\text{boot sample}) = \left(1 - \frac{1}{n}\right)^n.$$

\medskip
\sol{} \subpoints{1} \textbf{(ii)} Standard limit:
$$\lim_{n\to\infty}\left(1-\frac{1}{n}\right)^n = e^{-1} \approx \boxed{0.368},$$
so $P(i\in\text{boot sample})\to 1 - e^{-1} \approx \boxed{0.632}$ for large $n$.

\medskip
\sol{} \subpoints{1} \textbf{(iii)} A bootstrap estimate of $\mathrm{SE}(\hat\theta)$ for the sample median:
\begin{enumerate}
  \item Set the number of bootstrap resamples $B$ (e.g.\ $B=1000$).
  \item For $b=1,\dots,B$: draw $(X_1^{*b},\dots,X_n^{*b})$ with replacement from the original sample, and
        compute $\hat\theta^{*b} = \mathrm{median}(X_1^{*b},\dots,X_n^{*b})$.
  \item Form the bootstrap mean $\bar\theta^* = \frac{1}{B}\sum_{b=1}^B\hat\theta^{*b}$.
  \item Return
        $$\widehat{\mathrm{SE}}_{\text{boot}}(\hat\theta) =
          \sqrt{\frac{1}{B-1}\sum_{b=1}^B\left(\hat\theta^{*b} - \bar\theta^*\right)^2}.$$
\end{enumerate}

\noindent\emph{Grading: full credit for any algorithm that draws with replacement, recomputes $\hat\theta$ on
each resample, and returns the sample SD of the bootstrap replicates. Deduct 0.5 if sampling without
replacement is used.}

\vspace{0.6em}
\hrule
\vspace{1em}

%=========================================================
\section*{Problem 4 \points{22} --- Data analysis: regression on house prices}
%=========================================================

\subsection*{a) Ordinary least squares with an interaction \subpoints{5}}

\sol{} \subpoints{1} \textbf{(i)} Count the rows of the coefficient table: intercept, \texttt{crim},
\texttt{rm}, \texttt{age}, \texttt{dis}, \texttt{tax}, \texttt{ptratio}, \texttt{chas}, \texttt{rm:chas}
$=\boxed{9}$ parameters. Consistency check: residual d.f.\ $= n_\text{train} - p = 400 - 9 = 391$, matching
the printed residual d.f.\ exactly. \checkmark

\medskip
\sol{} \subpoints{2} \textbf{(ii)} The classmate misreads the table. In a model with an interaction term
\texttt{rm:chas}, the coefficient on \texttt{chas} alone is the river effect \emph{at} \texttt{rm}\,$=0$
(standardized, i.e.\ at the mean number of rooms); the full effect of being next to the river is
$\hat\beta_{\text{chas}} + \hat\beta_{\text{rm:chas}}\cdot\texttt{rm}$, which depends on \texttt{rm}.
A non-significant main effect on \texttt{chas} does not imply that \texttt{chas} is irrelevant globally;
the interaction is significant ($p=0.027$) and pulls the river effect upward for high-\texttt{rm} tracts.

Implied river effects (in $1000$\,USD):
\begin{itemize}
  \item (a) at standardized \texttt{rm}\,$=0$ (the mean):
        $\hat\beta_{\text{chas}} + \hat\beta_{\text{rm:chas}}\cdot 0 = 1.20 + 0 = \boxed{1.20}.$
  \item (b) at standardized \texttt{rm}\,$=+1$ (one SD above mean):
        $\hat\beta_{\text{chas}} + \hat\beta_{\text{rm:chas}}\cdot 1 = 1.20 + 2.00 = \boxed{3.20}.$
\end{itemize}

\medskip
\sol{} \subpoints{1} \textbf{(iii)} Because the predictors are standardized, the coefficient
$\hat\beta_{\text{ptratio}} = -1.80$ is the model's predicted change in \texttt{price} (in $1000$ USD)
for a one-standard-deviation increase in the pupil--teacher ratio, holding the other predictors fixed.
A $95\%$ Wald confidence interval is
$$\hat\beta_{\text{ptratio}} \;\pm\; 1.96\cdot\widehat{\mathrm{SE}}\bigl(\hat\beta_{\text{ptratio}}\bigr)
   \;=\; -1.80 \;\pm\; 1.96\cdot 0.30 \;=\; -1.80 \pm 0.588
   \;=\; \boxed{[\,-2.39,\;-1.21\,]}\text{ (in $1000$ USD).}$$
So a one-SD rise in the pupil--teacher ratio is associated with a price drop of about $1{,}800$ USD on
average; the interval excludes zero, consistent with the highly significant $t$-value of $-6.00$ in the
table.

\noindent\emph{Grading: 1\,P for the correct CI (accept any minor rounding, e.g.\ $[-2.39,-1.21]$ or
$[-2.388,-1.212]$). Deduct $0.5$ if the standardization rationale is omitted.}

\medskip
\sol{} \subpoints{1} \textbf{(iv)} A single marginal $p$-value reflects \texttt{age}'s contribution
\emph{conditional on the other predictors}. \texttt{age} may be correlated with $\texttt{dis}$, $\texttt{tax}$,
or other predictors (e.g.\ old neighbourhoods are clustered, and proportion of pre-1940 housing is
collinear with distance to employment centres), and its signal can be soaked up. Dropping it on the basis
of one $t$-test alone is inappropriate; one should check whether the cross-validated test error rises
after removal, or compare predictive performance with and without the term, before dropping it.

\subsection*{b) Forward stepwise vs.\ best subset \subpoints{3}}

\sol{} \subpoints{1} \textbf{(i)} Forward stepwise fits the null model plus, at each step
$k=0,1,2,3$, the $7-k$ candidates that add one predictor on top of the current best size-$k$ model.
Through step $4$ (i.e.\ choosing the best $4$-predictor model) the total is
$$1 + 7 + 6 + 5 + 4 = \boxed{23} \text{ distinct candidate models.}$$

\medskip
\sol{} \subpoints{1} \textbf{(ii)} Best-subset fits every subset up to and including size $4$:
$$\binom{7}{0} + \binom{7}{1} + \binom{7}{2} + \binom{7}{3} + \binom{7}{4}
   = 1 + 7 + 21 + 35 + 35 = \boxed{99} \text{ candidate models.}$$

\medskip
\sol{} \subpoints{1} \textbf{(iii)} Forward stepwise is \emph{greedy}: each step locks in a predictor and
never revisits the choice. The best $4$-predictor model under best subset is the global RSS-minimiser over
all $\binom{7}{4}=35$ size-$4$ subsets, while the forward path is constrained to extend the
chosen $3$-predictor subset. When the optimal $4$-subset cannot be reached by sequentially adding to the
optimal $3$-subset --- a classic outcome with correlated predictors --- the two methods disagree.

\subsection*{c) Ridge with $10$-fold cross-validation \subpoints{7}}

\sol{} \subpoints{1} \textbf{(i)} Ridge minimises (with the intercept \emph{not} penalised, since standardized
predictors have mean zero so $\hat\beta_0 = \bar y$ regardless):
$$\hat\beta^R_\lambda \;=\; \arg\min_{\beta_0,\beta}\,
   \left\{\sum_{i=1}^{n}\bigl(y_i - \beta_0 - x_i^\top\beta\bigr)^2 \;+\; \lambda\sum_{j=1}^{p}\beta_j^2\right\},
   \qquad \lambda\geq 0.$$
The penalty sum runs over $j=1,\dots,p$ (here $p=7$), \emph{excluding} the intercept. Closed-form:
$\hat\beta^R_\lambda = (X^\top X + \lambda I_p)^{-1}X^\top y$ on centred data.

\medskip
\sol{} \subpoints{2} \textbf{(ii)} The ridge penalty $\lambda\sum_j\beta_j^2$ is smooth and strictly convex in
$\beta$. Its gradient is zero only at $\beta=0$ jointly; for finite $\lambda$ the minimiser shrinks each
$\hat\beta_j$ \emph{toward} zero but generically does not set any of them \emph{exactly} to zero (it would
require the unpenalised gradient to perfectly cancel the linear ridge derivative at exactly $\beta_j=0$,
a measure-zero event). Hence all seven coefficients remain nonzero at both $\hat\lambda_{\min}$ and
$\hat\lambda_{1\text{SE}}$.

The \emph{lasso} penalty $\lambda\sum_j|\beta_j|$ is non-differentiable at zero, and its subdifferential
contains an interval $[-\lambda,\lambda]$ at $\beta_j=0$. Whenever the gradient of the data term
$|2x_j^\top(y - X\beta)|$ is smaller than $\lambda$, the optimum is at $\beta_j=0$ exactly. So lasso
\emph{does} produce sparse coefficient vectors and would zero out one or more of the seven predictors at
the lasso CV-optimal $\lambda$.

\medskip
\sol{} \subpoints{1} \textbf{(iii)} The ridge penalty $\lambda\sum_j\beta_j^2$ treats the coefficients as
directly comparable. Without standardization, a predictor measured on a large numeric scale has
a tiny OLS coefficient and is penalised lightly, while one on a small scale has a large OLS coefficient
and is penalised heavily. Standardising every predictor to mean $0$, SD $1$ makes the penalty scale-invariant
in the units of the data and gives every predictor a fair chance of being shrunk.

\medskip
\sol{} \subpoints{2} \textbf{(iv)} Test MSE drops from $19.40$ (OLS) to $16.20$ (ridge at $\hat\lambda_{\min}$),
a $16.5\%$ reduction. Since ridge \emph{adds bias} relative to OLS but reduces variance, the only way ridge
can beat OLS in test MSE is if OLS is \emph{variance-dominated} on this dataset. With $n_\text{train}=400$
and $p=7$ standardized predictors, OLS is mostly fine but is sufficiently variance-heavy ---
likely due to correlation among the continuous predictors (\texttt{age}, \texttt{dis}, \texttt{tax} are
correlated, for instance) --- that shrinking coefficients toward zero buys more in variance reduction than
it pays in bias inflation. Ridge is winning the bias--variance trade-off here.

\medskip
\sol{} \subpoints{1} \textbf{(v)} \textbf{Comparable or slightly lower than OLS.} If most predictors carry
real signal (which the OLS $t$-table suggests --- only \texttt{age} and \texttt{chas} have $p>0.05$), the
lasso will zero out at most one or two coefficients and shrink the rest, with a test MSE roughly between
the OLS value of $19.40$ and the ridge value of $16.20$. A defensible alternative answer: ``slightly higher
than ridge but lower than OLS, since the data has weak sparsity.''

\subsection*{d) PCR and spline degrees of freedom \subpoints{4}}

\sol{} \subpoints{2} \textbf{(i)} Total variance (sum of eigenvalues of the correlation matrix):
$$\sum_{j=1}^7\lambda_j = 2.80 + 1.60 + 1.00 + 0.70 + 0.40 + 0.30 + 0.20 = 7.00.$$
(Sanity check: standardized predictors have variance $1$, so the trace of the correlation matrix is exactly
$p=7$. \checkmark) The $80\%$ threshold is $0.80\cdot 7.00 = 5.60$. Cumulative sums:
\begin{align*}
\text{PC1:}    &\quad 2.80 / 7.00 = 40.00\% \\
\text{PC1+2:}  &\quad 4.40 / 7.00 = 62.86\% \\
\text{PC1+3:}  &\quad 5.40 / 7.00 = 77.14\%\;\; (<80\%) \\
\text{PC1+4:}  &\quad 6.10 / 7.00 = 87.14\%\;\; (\geq 80\%) \;\;\checkmark
\end{align*}
Need $\boxed{4}$ principal components to explain at least $80\%$ of the total variance.

\medskip
\sol{} \subpoints{1} \textbf{(ii)} Both PCR and ridge place \emph{less weight on low-variance directions} of
the predictor space: write the standardized design matrix as $X = U D V^\top$ (SVD). PCR keeps the top $M$
principal components --- it gives full weight to the first $M$ directions and \emph{zero} weight to the rest.
That is the ``discrete'' aspect: directions are either fully in or fully out. Ridge instead applies the
weight $d_j^2/(d_j^2+\lambda)$ to the $j$th SVD direction, which is a smooth, monotone shrinkage from
$1$ (large $d_j$, large variance) down toward $0$ (small $d_j$, small variance). Same shape of pressure,
continuous (ridge) versus discrete (PCR).

\medskip
\sol{} \subpoints{1} \textbf{(iii)} Counting d.f.\ including the intercept:
\begin{itemize}
  \item Intercept: $1$ d.f.
  \item Cubic spline on \texttt{rm} with $K=3$ interior knots: $K+3=6$ d.f.\ (excluding intercept).
  \item Five remaining continuous predictors \texttt{crim}, \texttt{age}, \texttt{dis}, \texttt{tax},
        \texttt{ptratio} entered linearly: $5\cdot 1 = 5$ d.f.
  \item Binary \texttt{chas} (one dummy): $1$ d.f.
\end{itemize}
Total:
$$1 + 6 + 5 + 1 \;=\; \boxed{13}\text{ d.f.}$$

\subsection*{e) Boosting interpretation \subpoints{3}}

\sol{} \subpoints{1} \textbf{(i)} Boosting at test MSE $11.40$ versus OLS at $19.40$ and ridge at $16.20$
suggests the relationship between predictors and \texttt{price} contains substantial
\emph{nonlinearities and/or interactions} that the linear models cannot capture. Depth-$4$ trees can model
up to $4$-way interactions and the smooth (often nonlinear) effect of \texttt{rm} on price.

\medskip
\sol{} \subpoints{1} \textbf{(ii)} Not a good idea. Unlike random forests, gradient boosting \emph{can}
overfit for very large $B$: each new tree fits residuals of the running ensemble and, with enough trees,
those residuals are just noise. $B$ is a genuine tuning parameter; ``$B=10{,}000$ just to be safe'' will
hurt test MSE. The right procedure is to monitor CV / OOB / validation error as $B$ grows and stop when
it stops decreasing.

\medskip
\sol{} \subpoints{1} \textbf{(iii)} Two diagnostics for a boosted ensemble:
\begin{itemize}
  \item \textbf{Variable importance plot} (mean decrease in node impurity or permutation importance).
        \emph{Answers}: which predictors the ensemble uses most heavily. \emph{Does not answer}: the
        \emph{direction} or shape of each predictor's effect, or which interactions are present.
  \item \textbf{Partial dependence plot} on a single predictor.
        \emph{Answers}: the marginal shape of how the predicted \texttt{price} depends on that predictor,
        averaged over the others. \emph{Does not answer}: how the effect changes within subgroups (i.e.\
        interactions, unless you compute 2-D PDPs) or whether the marginal effect is causal.
\end{itemize}

\vspace{0.6em}
\hrule
\vspace{1em}

%=========================================================
\section*{Problem 5 \points{24} --- Data analysis: classification of heart disease}
%=========================================================

\subsection*{a) Discriminant analysis: LDA vs.\ QDA \subpoints{7}}

\sol{} \subpoints{2} \textbf{(i)} With shared $\hat\sigma^2 = 2$ and equal priors $\hat\pi_0=\hat\pi_1=0.5$
(so $\log\hat\pi_0=\log\hat\pi_1=\log 0.5 = -\log 2$):
\begin{align*}
\delta_0(x) &= x\cdot\frac{\hat\mu_0}{\hat\sigma^2} - \frac{\hat\mu_0^2}{2\hat\sigma^2} + \log\hat\pi_0
   = x\cdot\frac{4}{2} - \frac{16}{4} + \log(0.5) \\
   &= 2x - 4 + \log(0.5), \\[4pt]
\delta_1(x) &= x\cdot\frac{\hat\mu_1}{\hat\sigma^2} - \frac{\hat\mu_1^2}{2\hat\sigma^2} + \log\hat\pi_1
   = x\cdot\frac{6}{2} - \frac{36}{4} + \log(0.5) \\
   &= 3x - 9 + \log(0.5).
\end{align*}

\medskip
\sol{} \subpoints{1} \textbf{(ii)} Set $\delta_0(x)=\delta_1(x)$. The shared $\log(0.5)$ cancels:
\begin{align*}
2x - 4 &= 3x - 9 \\
9 - 4 &= 3x - 2x \\
5 &= x.
\end{align*}
Decision boundary at $\boxed{x=5}$. Interpretation: with equal priors and equal variances, the LDA boundary
sits at the \emph{midpoint} $(\hat\mu_0 + \hat\mu_1)/2 = (4+6)/2 = 5$ of the two class means; below $5$ we
predict class $0$ (no CHD), above $5$ we predict class $1$ (CHD).

\medskip
\sol{} \subpoints{1} \textbf{(iii)} With $\hat\pi_0=0.8$, $\hat\pi_1=0.2$, the priors no longer cancel.
Solving $\delta_0(x)=\delta_1(x)$:
\begin{align*}
2x - 4 + \log(0.8) &= 3x - 9 + \log(0.2) \\
2x - 3x &= -9 + \log(0.2) - (-4 + \log(0.8)) \\
-x &= -5 + \log(0.2) - \log(0.8) \\
-x &= -5 + \log\!\left(\tfrac{0.2}{0.8}\right) = -5 + \log(0.25) = -5 - \log 4 \\
x &= 5 + \log 4 \approx 5 + 1.386 = \boxed{6.386}.
\end{align*}
The boundary moves \emph{toward} the class-$1$ mean ($\hat\mu_1=6$). This makes sense: with a strong prior
$\hat\pi_0=0.8$ on ``no CHD,'' we need more evidence (a higher $x$) before classifying as CHD, so the
threshold moves rightward.

\medskip
\sol{} \subpoints{3} \textbf{(iv)} With $\hat\sigma_0^2=1$, $\hat\sigma_1^2=4$, $\hat\mu_0=4$, $\hat\mu_1=6$,
$\hat\pi_0=\hat\pi_1=0.5$:
\begin{align*}
\delta_0(x) &= -\frac{x^2}{2\hat\sigma_0^2} + \frac{x\hat\mu_0}{\hat\sigma_0^2}
              - \frac{\hat\mu_0^2}{2\hat\sigma_0^2} - \tfrac{1}{2}\log\hat\sigma_0^2 + \log\hat\pi_0 \\[2pt]
            &= -\frac{x^2}{2\cdot 1} + \frac{4x}{1} - \frac{16}{2\cdot 1} - \tfrac{1}{2}\log 1 + \log(0.5) \\[2pt]
            &= -\tfrac{1}{2}x^2 + 4x - 8 + 0 + \log(0.5), \\[6pt]
\delta_1(x) &= -\frac{x^2}{2\hat\sigma_1^2} + \frac{x\hat\mu_1}{\hat\sigma_1^2}
              - \frac{\hat\mu_1^2}{2\hat\sigma_1^2} - \tfrac{1}{2}\log\hat\sigma_1^2 + \log\hat\pi_1 \\[2pt]
            &= -\frac{x^2}{2\cdot 4} + \frac{6x}{4} - \frac{36}{2\cdot 4} - \tfrac{1}{2}\log 4 + \log(0.5) \\[2pt]
            &= -\tfrac{1}{8}x^2 + \tfrac{3}{2}x - \tfrac{9}{2} - \tfrac{1}{2}\log 4 + \log(0.5).
\end{align*}

The QDA boundary is quadratic because, when class variances differ, the $-x^2/(2\sigma_k^2)$ term has a
different coefficient in $\delta_0$ and $\delta_1$. Subtracting them leaves an $x^2$ term that does \emph{not}
cancel:
$$\delta_0(x) - \delta_1(x) = \left(-\tfrac{1}{2} + \tfrac{1}{8}\right)x^2 + \cdots
   = -\tfrac{3}{8}x^2 + \cdots,$$
so the boundary $\delta_0=\delta_1$ is a quadratic equation in $x$ (generically two roots, hence a
quadratic decision region). Under the LDA assumption $\sigma_0^2=\sigma_1^2$, the $x^2$ coefficients are
equal and cancel, leaving a linear boundary.

\noindent\emph{Grading: 1\,P each correct discriminant; 1\,P for the structural explanation (the $x^2$
coefficients differ when variances differ). Deduct 1\,P per algebra mistake.}

\subsection*{b) Logistic regression with an interaction \subpoints{8}}

\sol{} \subpoints{2} \textbf{(i)} Encoding: $\texttt{sex}=0$ for male (reference), $\texttt{sex}=1$ for female.

For one extra year of \texttt{age}, holding everything else fixed, the log-odds change is
$\hat\beta_{\text{age}} + \hat\beta_{\text{age:sex}}\cdot\texttt{sex}$:
\begin{itemize}
  \item (a) Male ($\texttt{sex}=0$): change in log-odds $=\hat\beta_{\text{age}} = 0.06$,
        so odds factor $= e^{0.06} \approx \boxed{1.06}$.
  \item (b) Female ($\texttt{sex}=1$): change in log-odds
        $=\hat\beta_{\text{age}} + \hat\beta_{\text{age:sex}} = 0.06 + (-0.04) = 0.02$,
        so odds factor $= e^{0.02} \approx \boxed{1.02}$.
\end{itemize}

\noindent\emph{Grading: 1\,P male, 1\,P female. Deduct 1 if both factors are $e^{0.06}$, which ignores
the interaction.}

\medskip
\sol{} \subpoints{1} \textbf{(ii)} For one extra mmol/L of \texttt{ldl}, the odds change by
$$e^{\hat\beta_{\text{ldl}}} = e^{0.40} \approx \boxed{1.49}.$$
This factor is the same for males and females because the model has \emph{no} $\texttt{ldl}\!:\!\texttt{sex}$
interaction: the slope on $\texttt{ldl}$ on the log-odds scale is a single number $\hat\beta_1=0.40$ shared
across both sex strata. (Only \texttt{age} interacts with \texttt{sex} in this model.)

\medskip
\sol{} \subpoints{2} \textbf{(iii)} The patient: male ($\texttt{sex}=0$), $\texttt{age}=55$,
$\texttt{ldl}=5$, $\texttt{famhist}=1$. The linear predictor is
\begin{align*}
\hat\eta &= \hat\beta_0 + \hat\beta_1\,\texttt{ldl} + \hat\beta_2\,\texttt{age}
          + \hat\beta_3\,\texttt{sex} + \hat\beta_4\,(\texttt{age}\!:\!\texttt{sex})
          + \hat\beta_5\,\texttt{famhist} \\
   &= -5.00 + 0.40\cdot 5 + 0.06\cdot 55 + 1.20\cdot 0 + (-0.04)\cdot(55\cdot 0) + 0.80\cdot 1 \\
   &= -5.00 + 2.00 + 3.30 + 0 + 0 + 0.80 \\
   &= \boxed{1.10}.
\end{align*}
Fitted probability:
$$\hat p = \frac{e^{\hat\eta}}{1 + e^{\hat\eta}}
        = \frac{e^{1.10}}{1 + e^{1.10}}
        = \frac{3.0042}{4.0042}
        \approx \boxed{0.7503}.$$
At threshold $0.5$, $\hat p > 0.5$, so the predicted class is $\boxed{\hat y = 1}$ (predicted CHD).

\noindent\emph{Grading: 1\,P linear predictor; 1\,P probability + class label. Accept $\hat p\in[0.74,0.76]$
for rounding.}

\medskip
\sol{} \subpoints{1} \textbf{(iv)} When a model contains an interaction $\texttt{age}\!:\!\texttt{sex}$, the
main-effect coefficient $\hat\beta_{\text{sex}}=1.20$ is the sex effect \emph{at} $\texttt{age}=0$, which is
far outside the observed data range. Its $p$-value tests whether the sex effect is zero \emph{at age zero},
which is essentially meaningless here. Dropping the main-effect \texttt{sex} term while keeping the
interaction would force the sex effect to be zero at $\texttt{age}=0$, distorting the fitted effects at
realistic ages. The standard rule (``hierarchy / marginality principle'') is: do not drop a main effect
whose interaction is in the model.

\medskip
\sol{} \subpoints{2} \textbf{(v)} In the CHD setting:
\begin{itemize}
  \item \textbf{Sensitivity} = $P(\hat Y = 1 \mid Y = 1)$: the proportion of patients who \emph{actually
        develop CHD} that the classifier correctly flags as CHD. (``Of the patients who really get sick,
        how many do we catch?'')
  \item \textbf{Specificity} = $P(\hat Y = 0 \mid Y = 0)$: the proportion of patients who \emph{never
        develop CHD} that the classifier correctly clears. (``Of the healthy patients, how many do we
        leave alone?'')
\end{itemize}

\subsection*{c) ROC curve and threshold tuning \subpoints{5}}

The test set has $n=200$ patients: $80$ actual CHD (TP $=30$, FN $=50$) and $120$ actual non-CHD
(TN $=104$, FP $=16$).

\medskip
\sol{} \subpoints{2} \textbf{(i)}
\begin{align*}
\text{Sensitivity} &= \frac{\text{TP}}{\text{TP}+\text{FN}}
   = \frac{30}{30+50} = \frac{30}{80} = \boxed{0.375}, \\[3pt]
\text{Specificity} &= \frac{\text{TN}}{\text{TN}+\text{FP}}
   = \frac{104}{104+16} = \frac{104}{120} \approx \boxed{0.867}, \\[3pt]
\text{Test error rate} &= \frac{\text{FP}+\text{FN}}{n}
   = \frac{16+50}{200} = \frac{66}{200} = \boxed{0.33}.
\end{align*}

\noindent\emph{Grading: 1\,P for sensitivity, 1\,P for specificity + error.}

\medskip
\sol{} \subpoints{1} \textbf{(ii)} The \textbf{AUC} is the area under the ROC curve, swept over all
classification thresholds in $[0,1]$. Equivalently, AUC equals the probability that the classifier ranks
a randomly chosen positive observation higher (by predicted probability) than a randomly chosen negative
observation. $\boxed{\mathrm{AUC}=0.5}$ corresponds to a classifier that ranks no better than random
guessing (the diagonal line in ROC space).

\medskip
\sol{} \subpoints{2} \textbf{(iii)} Lowering the threshold from $0.5$ to $0.3$ flags \emph{more} patients
as CHD-positive:
\begin{enumerate}[(a)]
  \item Sensitivity \emph{increases} (more true positives are caught --- jumping from $0.375$ to roughly
        $0.66$ per the marked point on the ROC curve), and specificity \emph{decreases} (more false
        alarms, jumping from $\approx 0.867$ to $\approx 0.64$, i.e.\ FPR rises from $\sim 0.13$ to $\sim 0.36$).
  \item On the ROC curve the operating point moves \emph{up and to the right}, from the marked $\hat p=0.5$
        point to the marked $\hat p=0.3$ point.
\end{enumerate}
A clinical context where this is the right move: \emph{screening} for CHD in a setting where missing a true
case is much more costly (delayed treatment, future events) than a false alarm, which only triggers a
follow-up test.

\subsection*{d) A tree-based competitor \subpoints{4}}

\sol{} \subpoints{2} \textbf{(i)} I would fit a \textbf{random forest} for binary classification.
Tuning parameters and their justifications:
\begin{itemize}
  \item $B$ (number of trees) $= 500$. Random forests do not overfit in $B$; we keep growing until
        the OOB / test error stabilises, and $500$ is the conventional plateau point.
  \item $m$ (predictors sampled per split) $= \lfloor\sqrt{p}\rfloor$. With $p=6$ candidate
        predictors that gives $m=2$ (or $3$). The role of $m$ is to \emph{decorrelate} the trees
        (lower $\rho$ in the bagging variance formula from P3a) and reduce variance.
  \item Minimum terminal node size $=1$ (i.e.\ trees grown deep). Variance is controlled by
        averaging across the $B$ trees rather than by per-tree pruning; deep trees give low-bias
        base learners.
\end{itemize}

\noindent\emph{Grading: 2\,P for naming the method and giving a concrete value for each of three
parameters with one-sentence justification. Deduct 0.5 per missing parameter or missing justification.}

\medskip
\sol{} \subpoints{1} \textbf{(ii)} Forest: sensitivity $0.50$, specificity $0.85$, error rate $0.21$.
Logistic at threshold $0.5$: sensitivity $0.375$, specificity $0.867$, error $0.33$. The random forest
\emph{wins} on test error rate ($0.21<0.33$) and \emph{wins} on sensitivity ($0.50>0.375$), at a small
cost in specificity ($0.85<0.867$). In a clinical CHD-screening setting where catching true cases matters
more than avoiding the occasional false alarm, the higher sensitivity and lower overall error rate make
\textbf{random forest} the recommended classifier. The decisive criterion is the combination of
test error and sensitivity in a context where false negatives are clinically costly.

\medskip
\sol{} \subpoints{1} \textbf{(iii)} With $\Pr(\texttt{chd}=1)\approx 0.27$, the trivial ``always
predict no CHD'' classifier already attains test error $\approx 0.27$ --- which actually
\emph{beats} the logistic regression's error of $0.33$. Test error alone is therefore a misleading
metric: a useless classifier can ``win'' it on imbalanced data. The right metrics here are the
\emph{(sensitivity, specificity)} pair (which separately measure performance on each class) and the
threshold-independent \textbf{AUC}; in a cost-sensitive screening context, one should privilege
sensitivity (or the $F_1$ score, which weighs precision and recall) over raw error rate.

\vspace{1em}
\hrule
\vspace{0.5em}

\noindent\textbf{End of solution proposal.}\quad Total awarded: $10+28+16+22+24=100$ points.

\end{document}
